\section{Evaluation}

\subsection{Clean Test Results}
\begin{table}[H]
  \centering
  \begin{tabular}{@{}lrrr@{}}
    \toprule
    Metric & Baseline & FP32 Edge Model & Quantized Model \\
    \midrule
    Accuracy         &  &  &  \\
    Macro F1         &  &  &  \\
    Model parameters &  &  &  \\
    \bottomrule
  \end{tabular}
  \caption{Clean DS2 test results.}
  \label{tab:clean-results}
\end{table}

\subsection{Per-Class Performance}
\begin{table}[H]
  \centering
  \begin{tabular}{@{}lrrr@{}}
    \toprule
    Class & Precision & Sensitivity/Recall & F1 \\
    \midrule
    N &  &  &  \\
    S &  &  &  \\
    V &  &  &  \\
    F &  &  &  \\
    \bottomrule
  \end{tabular}
  \caption{Per-class performance of the final model on DS2.}
  \label{tab:per-class}
\end{table}
% Confusion matrix figure; discuss which classes are difficult and why.

\subsection{Quantization Analysis}
\begin{itemize}
  \item FP32 Macro F1: \rule{2cm}{0.4pt}
  \item INT8 Macro F1: \rule{2cm}{0.4pt}
  \item Difference: \rule{2cm}{0.4pt}
\end{itemize}

\subsection{Noise Robustness}
\begin{table}[H]
  \centering
  \begin{tabular}{@{}lrr@{}}
    \toprule
    Condition & Accuracy & Macro F1 \\
    \midrule
    Clean &  &  \\
    12 dB &  &  \\
    6 dB  &  &  \\
    0 dB  &  &  \\
    \bottomrule
  \end{tabular}
  \caption{Performance under increasing ECG noise.}
  \label{tab:noise}
\end{table}
% Plot performance vs noise level. Discuss: most affected classes, gradual degradation,
% whether preprocessing helps, whether robustness could be improved.

\subsection{Error Analysis}
% Representative errors with ECG waveforms: N->S, S->N, V misclassified, errors under strong noise.

\subsection{Success Criteria Review}
% For each criterion from Section 1.4: Achieved / Partially achieved / Not achieved, with justification.
